\documentclass[12pt]{article}

\usepackage[T2A]{fontenc}
\usepackage[utf8]{inputenc}
\usepackage[russian]{babel}
\usepackage{amsmath,amssymb}
\usepackage{microtype}
\usepackage{fullpage}
\usepackage{multicol,multirow}
\usepackage{tabularx}

% Оригиналный шаблон: http://k806.ru/dalabs/da-report-template-2012.tex

\begin{document}

\section*{Курсовая работа по курсу дискретного анализа: Алгоритм LZ77}

Выполнил студент группы 08-303 МАИ \textit{Арусланов Кирилл}.

\subsection*{Условие}

Реализовать алгоритм LZ77. По условию требуется сжимать и восстанавливать текст, состоящий только из малых латинских букв. Поиск совпадений производится не по скользящему окну, а по всему уже просмотренному ранее тексту. Формат ввода/вывода:
\begin{itemize}
  \item Команда \texttt{compress} и далее текст --- вывести последовательность троек $\langle\text{offset},\,\text{len},\,\text{char}\rangle$, по одной тройке на строку. Если на последней тройке символ отсутствует, не выводить его.
  \item Команда \texttt{decompress} и далее тройки --- восстановить исходный текст. Тройка может не содержать символ (это означает конец текста).
\end{itemize}

\subsection*{Метод решения}

В работе реализован классический подход поиска наибольшего совпадения с ранее просмотренной частью строки, но для эффективного поиска используется суффиксный массив (Suffix Array) и массив длин общих префиксов (LCP). Основные идеи:
\begin{enumerate}
  \item Строится преобразованная строка с явным терминатором минимального кода (функция \texttt{map\_with\_terminator}). Это делает все суффиксы сравнимыми и корректно обрабатывает границы, также сдвиг способствует уменьшению затрат по памяти при сортировке подсчетом.
  \item Для строки строится суффиксный массив методом сортировки по классам ($O(n\log n)$ в практике) --- функция \texttt{build\_sa\_mapped}.
  \item На основании суффиксного массива рассчитывается LCP (массив длин общих префиксов) --- функция \texttt{build\_lcp\_mapped}. Это даёт возможность быстро оценивать длину совпадения между текущей позицией и любой позицией в предыдущей части текста.
  \item Для каждой позиции $pos$ берутся соседи по SA (в обе стороны) и по массиву LCP ищется наибольшее вхождение, которое начинается строго раньше $pos$ (т.е. в ранее просмотренной части). Если наибольшая длина равна нулю, кодируется одиночный символ как $0,0,c$; иначе кодируется тройка $offset,poslen,nextchar$ или, если совпадение выходит до конца текста, тройка без символа.
\end{enumerate}

Сложности по времени и памяти:
\begin{itemize}
  \item Построение SA — около $O(n\log n)$ по времени и $O(n)$ по памяти (включая строки и массивы). Для практических размеров (до нескольких сотен тысяч символов) работает эффективно.
  \item Вычисление LCP — $O(n)$.
  \item Проход по позициям и выбор лучшего вхождения — в худшем случае может быть до $O(n^2)$ при крайне неблагоприятной структуре и если поиск соседей долго просматривает, но на реальных данных благодаря LCP обычно быстрее.
\end{itemize}

\subsection*{Описание программы}

Программа реализована в одной точке входа (файл с кодом \texttt{main.cpp}). Основные функции/компоненты:
\begin{itemize}
  \item \texttt{map\_with\_terminator(const std::string\&)} --- сопоставляет символы строки смещёнными кодами и добавляет терминатор (байт 0), чтобы суффиксы сравнивались корректно.
  \item \texttt{build\_sa\_mapped(const std::string\&)} --- строит суффиксный массив методом сортировки по классам (алгоритм типа \textit{prefix doubling}). Возвращает вектор позиций суффиксов.
  \item \texttt{build\_lcp\_mapped(const std::string\&, const std::vector<int>\&sa)} --- строит LCP во времени $O(n)$ по алгоритму Касаи.
  \item \texttt{lz77\_compress\_sa(const std::string\&)} --- основная процедура сжатия: для каждой позиции находит лучшее совпадение в предыдущей части с помощью SA/LCP и формирует вектор троек (структура \texttt{Token}).
  \item \texttt{lz77\_decompress(const std::vector<Token>\&)} --- выполняет восстановление строки из набора троек, корректно поддерживает перекрывающиеся копии.
  \item \texttt{parse\_tokens\_from\_stream} (в улучшенной версии) --- парсит входные строки с тройками построчно, устойчив к лишним пробелам.
  \item \texttt{main} --- парсит команду \texttt{compress}/\texttt{decompress} в одном из двух форматов (команда и данные либо в одной строке, либо на следующей) и вызывает соответствующие функции.
\end{itemize}

\subsection*{Дневник отладки}

Работа была реализована и протестирована на наборе тестов, обеспечивающих корректность и граничные случаи. Краткая хронология и замечания:
\begin{itemize}
  \item Реализация основных алгоритмов SA и LCP. На этапе раннего тестирования проверялась корректность на простых строках: \texttt{"a"}, \texttt{"aaaaa"}, \texttt{"abababab"}.
  \item Доработана обработка пробелов при декодировании.
  \item Скорректированы границы цикла при кодировании.
\end{itemize}

\subsection*{Тест производительности}

Измерения времени работы производились для случайно сгенерированных строк различной длины. Время — wall-clock (мс), замеры одного прогона на каждой длине. В таблице приведено время сжатия, время распаковки и число полученных токенов.

\begin{center}
\begin{tabular}{r r r r}
\hline
$n$ & compress (ms) & decompress (ms) & tokens \\
\hline
1000 & 6 & 0 & 413 \\
5000 & 33 & 0 & 1698 \\
10000 & 51 & 0 & 3176 \\
30000 & 140 & 0 & 8577 \\
60000 & 179 & 0 & 16080 \\
120000 & 422 & 1 & 30466 \\
240000 & 939 & 4 & 58112 \\
480000 & 2301 & 8 & 110526 \\
\hline
\end{tabular}
\end{center}

\paragraph{Краткий разбор результатов.} По таблице видно, что:
\begin{itemize}
\item Время распаковки пренебрежимо мало по сравнению со временем сжатия и растёт очень медленно (близко к линейному поведению и с маленькой константой).
\item Время сжатия растёт существенно при увеличении $n$; это соответствует тому, что доминирующая часть работы — построение суффиксного массива и дополнительные шаги (теоретически близко к $O(n\log n)$ для использованного алгоритма), а также последующая обработка.
\item Число токенов растёт примерно линейно с $n$, при этом отношение токенов к длине строки уменьшается для больших $n$, что характерно для случайных данных (в относительном выражении вероятность длинных совпадений растёт с объёмом данных).
\end{itemize}

\subsection*{Выводы}

Реализованный алгоритм LZ77 с поиском по всему пройденному тексту и использованием суффиксного массива является корректным и пригоден для задач средней размерности. Применение:
\begin{itemize}
  \item Архиваторы/компоненты сжатия потока, где требуется найти повторяющиеся подстроки в данных.
  \item Анализ повторяющихся паттернов в текстах и последовательностях (биоинформатика, сжатие логов и т.п.).
\end{itemize}

Сложность реализации умеренная: основная сложность связана со стандартными структурами данных (SA/LCP).

\end{document}
